\section{Conclusion}
\label{sec:conclusion}

Our companion negative result left one question open: is a matrix latent's
rank ignored by the gradient in general, or only when nothing rewards it?
Closing every shortcut that could make a model look rank-aware without
needing rank -- argmax decoding, position-decomposition, and a step budget
too short to see a late transition -- the answer is that the gradient does
see rank, when a task provably forces it. At $d\in\{8,16\}$, gradient
descent trained from scratch recruits effective rank $\approx K$ and makes
it causally necessary, with a step at $k\approx K$ sharp enough to be
razor-clean at the smallest tested cell. That rank composes: a
rank-sufficient trained operator self-applies exactly to depth 21, and an
entity-subspace decomposition shows this is not a lucky lookup table but a
genuine, restricted-rank-exactly-$K$ realization of the target permutation,
with an invisible, unconstrained complement absorbing everything the loss
never asked it to control. Depth itself is a sharper rank probe than any
single-hop tolerance: a deliberately rank-starved operator's entire
depth-decay curve is predictable from its own spectrum, no raw keys
required. And what looked like a hard trainability wall at larger $d$ is
substantially a step-budget artifact -- three independent axes of this
program made the same mistake before a $2$--$2.5\times$ budget extension
corrected it -- with a genuine, narrower frontier surviving underneath: not
whether training moves, but how exactly it converges as $d$ grows. Every
number here traces to a specific run, seed, and step count, and every
claim above has a stated falsification condition. The transfer question --
whether any of this helps reasoning on real data -- is not addressed here
by design, and is the explicit next step this line of results gates.
